% Options for packages loaded elsewhere
% Options for packages loaded elsewhere
\PassOptionsToPackage{unicode}{hyperref}
\PassOptionsToPackage{hyphens}{url}
%
\documentclass[
  english,
  russian,
  12pt,
  a4paper,
  DIV=11,
  numbers=noendperiod]{scrreprt}
\usepackage{xcolor}
\usepackage{amsmath,amssymb}
\setcounter{secnumdepth}{5}
\usepackage{iftex}
\ifPDFTeX
  \usepackage[T1]{fontenc}
  \usepackage[utf8]{inputenc}
  \usepackage{textcomp} % provide euro and other symbols
\else % if luatex or xetex
  \usepackage{unicode-math} % this also loads fontspec
  \defaultfontfeatures{Scale=MatchLowercase}
  \defaultfontfeatures[\rmfamily]{Ligatures=TeX,Scale=1}
\fi
\usepackage{lmodern}
\ifPDFTeX\else
  % xetex/luatex font selection
\fi
% Use upquote if available, for straight quotes in verbatim environments
\IfFileExists{upquote.sty}{\usepackage{upquote}}{}
\IfFileExists{microtype.sty}{% use microtype if available
  \usepackage[]{microtype}
  \UseMicrotypeSet[protrusion]{basicmath} % disable protrusion for tt fonts
}{}
\usepackage{setspace}
% Make \paragraph and \subparagraph free-standing
\makeatletter
\ifx\paragraph\undefined\else
  \let\oldparagraph\paragraph
  \renewcommand{\paragraph}{
    \@ifstar
      \xxxParagraphStar
      \xxxParagraphNoStar
  }
  \newcommand{\xxxParagraphStar}[1]{\oldparagraph*{#1}\mbox{}}
  \newcommand{\xxxParagraphNoStar}[1]{\oldparagraph{#1}\mbox{}}
\fi
\ifx\subparagraph\undefined\else
  \let\oldsubparagraph\subparagraph
  \renewcommand{\subparagraph}{
    \@ifstar
      \xxxSubParagraphStar
      \xxxSubParagraphNoStar
  }
  \newcommand{\xxxSubParagraphStar}[1]{\oldsubparagraph*{#1}\mbox{}}
  \newcommand{\xxxSubParagraphNoStar}[1]{\oldsubparagraph{#1}\mbox{}}
\fi
\makeatother


\usepackage{longtable,booktabs,array}
\usepackage{calc} % for calculating minipage widths
% Correct order of tables after \paragraph or \subparagraph
\usepackage{etoolbox}
\makeatletter
\patchcmd\longtable{\par}{\if@noskipsec\mbox{}\fi\par}{}{}
\makeatother
% Allow footnotes in longtable head/foot
\IfFileExists{footnotehyper.sty}{\usepackage{footnotehyper}}{\usepackage{footnote}}
\makesavenoteenv{longtable}
\usepackage{graphicx}
\makeatletter
\newsavebox\pandoc@box
\newcommand*\pandocbounded[1]{% scales image to fit in text height/width
  \sbox\pandoc@box{#1}%
  \Gscale@div\@tempa{\textheight}{\dimexpr\ht\pandoc@box+\dp\pandoc@box\relax}%
  \Gscale@div\@tempb{\linewidth}{\wd\pandoc@box}%
  \ifdim\@tempb\p@<\@tempa\p@\let\@tempa\@tempb\fi% select the smaller of both
  \ifdim\@tempa\p@<\p@\scalebox{\@tempa}{\usebox\pandoc@box}%
  \else\usebox{\pandoc@box}%
  \fi%
}
% Set default figure placement to htbp
\def\fps@figure{htbp}
\makeatother



\ifLuaTeX
\usepackage[bidi=basic,provide=*]{babel}
\else
\usepackage[bidi=default,provide=*]{babel}
\fi
% get rid of language-specific shorthands (see #6817):
\let\LanguageShortHands\languageshorthands
\def\languageshorthands#1{}


\setlength{\emergencystretch}{3em} % prevent overfull lines

\providecommand{\tightlist}{%
  \setlength{\itemsep}{0pt}\setlength{\parskip}{0pt}}



 
\usepackage[style=gost-numeric,backend=biber,langhook=extras,autolang=other*]{biblatex}
\addbibresource{bib/cite.bib}

\usepackage[]{csquotes}

\usepackage{indentfirst}
\usepackage{float}
\floatplacement{figure}{H}
\IfFileExists{plex-otf.sty}{
  %% Full TeXlive
  \usepackage[
    % math,
    RM={Scale=0.94},SS={Scale=0.94},SScon={Scale=0.94},TT={Scale=MatchLowercase,FakeStretch=0.9},
    DefaultFeatures={Ligatures=Common}
  ]{plex-otf}
  % \usepackage[Scale=MatchUppercase]{juliamono}
}{
  %% TinyTeX
  \usepackage{libertine}
}
\KOMAoption{captions}{tableheading}
\makeatletter
\@ifpackageloaded{caption}{}{\usepackage{caption}}
\AtBeginDocument{%
\ifdefined\contentsname
  \renewcommand*\contentsname{Содержание}
\else
  \newcommand\contentsname{Содержание}
\fi
\ifdefined\listfigurename
  \renewcommand*\listfigurename{Список иллюстраций}
\else
  \newcommand\listfigurename{Список иллюстраций}
\fi
\ifdefined\listtablename
  \renewcommand*\listtablename{Список таблиц}
\else
  \newcommand\listtablename{Список таблиц}
\fi
\ifdefined\figurename
  \renewcommand*\figurename{Рисунок}
\else
  \newcommand\figurename{Рисунок}
\fi
\ifdefined\tablename
  \renewcommand*\tablename{Таблица}
\else
  \newcommand\tablename{Таблица}
\fi
}
\@ifpackageloaded{float}{}{\usepackage{float}}
\floatstyle{ruled}
\@ifundefined{c@chapter}{\newfloat{codelisting}{h}{lop}}{\newfloat{codelisting}{h}{lop}[chapter]}
\floatname{codelisting}{Список}
\newcommand*\listoflistings{\listof{codelisting}{Листинги}}
\makeatother
\makeatletter
\makeatother
\makeatletter
\@ifpackageloaded{caption}{}{\usepackage{caption}}
\@ifpackageloaded{subcaption}{}{\usepackage{subcaption}}
\makeatother
\usepackage{bookmark}
\IfFileExists{xurl.sty}{\usepackage{xurl}}{} % add URL line breaks if available
\urlstyle{same}
\hypersetup{
  pdftitle={Лабораторная работа №2},
  pdfauthor={Андрей Геннадьевич Софич},
  pdflang={ru-RU},
  hidelinks,
  pdfcreator={LaTeX via pandoc}}


\title{Лабораторная работа №2}
\usepackage{etoolbox}
\makeatletter
\providecommand{\subtitle}[1]{% add subtitle to \maketitle
  \apptocmd{\@title}{\par {\large #1 \par}}{}{}
}
\makeatother
\subtitle{Основные модели: SIR и модель Лотки-Вольтерры}
\author{Андрей Геннадьевич Софич}
\date{}
\begin{document}
\maketitle

\renewcommand*\contentsname{Содержание}
{
\setcounter{tocdepth}{1}
\tableofcontents
}
\listoffigures
\listoftables

\setstretch{1.5}
\chapter{Цель
работы}\label{ux446ux435ux43bux44c-ux440ux430ux431ux43eux442ux44b}

Познакомиться с базовым arduino uno

\chapter{Задание}\label{ux437ux430ux434ux430ux43dux438ux435}

В этом эксперименте мы научимся управлять светодиодом. Заставим его
мигать.

\chapter{Теоритическая
часть}\label{ux442ux435ux43eux440ux438ux442ux438ux447ux435ux441ux43aux430ux44f-ux447ux430ux441ux442ux44c}

Необходимые компоненты: •контроллер Arduino UNO R3; •плата для
прототипирования Breadboard; •светодиод; •резистор 220 Ом; •провода
Male-Male (папа-папа).

Светодиод -- это полупроводниковый прибор, преобразующий электрический
ток непосредственно в световое излучение. По-английски светодиод
называется light emitting diode, или LED. Цветовые характеристики
светодиодов зависят от химического состава использованного в нем
полупроводника. Светодиод излучает в узкой части спектра, его цвет чист,
что особенно ценят дизайнеры. Светодиод механически прочен и
исключительно надежен, его срок службы может достигать 100 тысяч часов,
что почти в 100 раз больше, чем у лампочки накаливания, и в 5--10 раз
больше, чем у люминесцентной лампы. Наконец, светодиод -- низковольтный
электроприбор, а стало быть, безопасный. Светодиоды поляризованы, имеет
значение, в каком направлении подключать их. Положительный вывод
светодиода (более длинный) называется анодом, отрицательный -- катодом.
Как и все диоды, светодиоды позволяют току течь только в одном
направлении -- от анода к катоду. Поскольку ток протекает от
положительного к отрицательному, анод светодиода должен быть подключен к
цифровому сигналу 5 В, а катод должен быть подключен к земле. Мы будем
подключать светодиод к цифровому контакту D10 Arduino последовательно с
резистором. Светодиоды должны быть всегда соединены последовательно с
резистором, который выступает в качестве ограничителя тока. Чем больше
значение резистора, тем больше он ограничивает ток. В этом эксперименте
мы используем резистор номиналом 220 Ом.

\chapter{Выполнение лабораторной
работы}\label{ux432ux44bux43fux43eux43bux43dux435ux43dux438ux435-ux43bux430ux431ux43eux440ux430ux442ux43eux440ux43dux43eux439-ux440ux430ux431ux43eux442ux44b}

Подключаем светодиод в плату для прототипирования (рис.~\ref{fig-001}).

\begin{figure}

\centering{

\includegraphics[width=0.7\linewidth,height=\textheight,keepaspectratio]{image/1.png}

}

\caption{\label{fig-001}Светодиоид}

\end{figure}%

Устанавливаем резистор, сопротивлением 220 ом в плату
(рис.~\ref{fig-002}).

\begin{figure}

\centering{

\includegraphics[width=0.7\linewidth,height=\textheight,keepaspectratio]{image/2.png}

}

\caption{\label{fig-002}Резистор}

\end{figure}%

Длинную ножку светодиода (анод) подключаем к цифровому выводу D10
Arduino, другую (катод) -- через резистор 220 Ом к выводу GND
(рис.~\ref{fig-003}).

\begin{figure}

\centering{

\includegraphics[width=0.7\linewidth,height=\textheight,keepaspectratio]{image/3.png}

}

\caption{\label{fig-003}Подключение}

\end{figure}%

Загружаем скетч в плату Arduino

const int LED=10; // вывод для подключения светодиода 10 (D10) void
setup() \{ // Конфигурируем вывод подключения светодиода как выход
(OUTPUT) pinMode(LED, OUTPUT); \} void loop() \{ // включаем светодиод,
подавая на вывод 1 (HIGH) digitalWrite(LED,HIGH); // пауза 1 сек (1000
мс) delay(1000); // выключаем светодиод, подавая на вывод 0 (LOW)
digitalWrite(LED,LOW); // пауза 1 сек (1000 мс) delay(1000); \}

Для мигания светодиода необходимо попеременно c определенным интервалом
подавать на вывод Arduino сигналы HIGH (высокий уровень или 1) и LOW
(низкий уровень или 0). Интервал изменения сигнала на выходе D10 Arduino
будем устанавливать с помощью команды delay(), задерживающей выполнение
скетча на заданное время в миллисекундах (мс).

Проверяем, получилось ли у нас успешно выполнить эксперимент.

\chapter{Выводы}\label{ux432ux44bux432ux43eux434ux44b}

В данной работе мы научились управлять светодиодом.


\printbibliography



\end{document}
